\section{Filtro}
En todo sistema de procesamiento de señales en tiempo discreto, el acondicionamiento analógico de la señal, tanto previo como posterior al proceso de discretización, constituye un pilar fundamental para garantizar la integridad de la información. De acuerdo con el Teorema de Muestreo de Nyquist-Shannon, existe un límite estricto sobre el ancho de banda máximo que una señal continua puede poseer para ser muestreada de manera unívoca. Superar este límite introduce componentes de alta frecuencia en la banda base, un fenómeno irreversible conocido como \textit{aliasing} o solapamiento espectral.

Por consiguiente, la implementación de etapas de filtrado analógico se vuelve imperativa. En la presente sección se abordará el análisis teórico, la viabilidad y las relaciones de compromiso de la red de filtrado activo requerida para este sistema. Se estudiará una arquitectura basada en celdas Sallen-Key, evaluando su desempeño en un rol dual: 
\begin{enumerate}
	\item Como \textbf{Filtro Anti-alias (FAA)}: encargado de acotar el espectro de la señal de entrada, limitando su ancho de banda antes de ingresar a la etapa de \textit{Sample and Hold}.
	\item Como \textbf{Filtro Recuperador (FR)}: posicionado a la salida del sistema para interpolar las muestras discretas, suprimiendo las imágenes espectrales de alta frecuencia (réplicas) generadas durante el muestreo y reconstruyendo así la señal continua original.
\end{enumerate}

Para satisfacer ambos requerimientos, la consigna propone la implementación de una topología activa pasa-bajos. El circuito propuesto, ilustrado en la \textbf{Figura \ref{fig:esquema_filtro_tp}}, se compone de amplificadores operacionales (TL082) configurados para formar dos celdas Sallen-Key en cascada seguidas de una red RC de primer orden. A continuación, se presenta dicho esquemático, el cual será sometido a un análisis matemático detallado para validar su orden y su respuesta en frecuencia frente a la plantilla exigida.
\begin{figure}[H]
	\centering
	\includegraphics[width=0.85\textwidth]{Imagenes/3_filtro/filtro.pdf} 
	\caption{Esquemático del filtro activo propuesto para las etapas FAA y FR.}
	\label{fig:esquema_filtro_tp}
\end{figure}
\newpage
\subsection{Filtro Anti-alias y Recuperador}
\subsubsection{Análisis de la Plantilla y el Filtro Propuesto}
Para determinar si la topología propuesta cumple con su cometido, se debe verificar matemáticamente si un filtro de orden 5 satisface la plantilla exigida. La consigna plantea un circuito basado en la aproximación de Chebyshev, implementado mediante dos celdas Sallen-Key y una etapa de primer orden.

Partimos de la plantilla pasa-bajos especificada:
\begin{itemize}
	\item Frecuencia de paso ($f_p$): $22.05\text{ kHz}$
	\item Frecuencia de atenuación ($f_a$): $44.1\text{ kHz}$
	\item Atenuación máxima en la banda de paso ($A_p$): $1.5\text{ dB}$
	\item Atenuación mínima en la banda de rechazo ($A_a$): $40\text{ dB}$
\end{itemize}

En primer lugar, calculamos el factor de rizado ($\epsilon$) característico de la banda de paso. Este parámetro se deduce de la función de magnitud de Chebyshev y sirve para cuantificar la amplitud de las ondulaciones permitidas antes de la frecuencia de corte:

\begin{equation}
	\epsilon = \sqrt{10^{A_p/10} - 1} = \sqrt{10^{0.15} - 1} \approx 0.642
\end{equation}

Luego, determinamos el orden mínimo ($n$) requerido para esta aproximación. El orden define físicamente la cantidad de polos del sistema y matemáticamente la pendiente de caída de la curva; a mayor orden, mayor es la selectividad. La inecuación se obtiene al igualar la función de atenuación de Chebyshev a los requisitos estrictos de la banda de rechazo:

\begin{equation}
	n \ge \frac{\cosh^{-1}\left(\frac{\sqrt{10^{A_a/10}-1}}{\epsilon}\right)}{\cosh^{-1}\left(\frac{f_a}{f_p}\right)}
\end{equation}

Sustituyendo los valores de la plantilla:

\begin{align*}
	n &\ge \frac{\cosh^{-1}\left(\frac{\sqrt{10^4-1}}{0.642}\right)}{\cosh^{-1}\left(\frac{44.1}{22.05}\right)} \\[1ex]
	n &\ge \frac{\cosh^{-1}(155.74)}{\cosh^{-1}(2)} \\[1ex]
	n &\ge \frac{5.74}{1.317} \approx 4.35
\end{align*}
\newpage
Dado que el orden $n$ debe ser estrictamente un número entero para poder sintetizarse físicamente, redondeamos al entero superior, obteniendo un requerimiento de $n = 5$. 

Por lo tanto, el filtro de orden 5 propuesto por la cátedra supera las restricciones de atenuación impuestas, confirmando su idoneidad tanto para la función de Filtro Anti-alias (acotando el espectro antes del muestreo para evitar aliasing) como de Filtro Recuperador (eliminando las réplicas espectrales post-muestreo).

\subsubsection{Consideraciones Prácticas de Implementación}

Si bien el diseño teórico de las celdas Sallen-Key garantiza el cumplimiento de la plantilla de atenuación, la implementación física del circuito exige la adición de componentes auxiliares para asegurar la estabilidad y la integridad de la señal frente a condiciones reales de operación. Por este motivo, el esquemático original ha sido modificado para incluir etapas de aislamiento y redes de filtrado en la alimentación.

\paragraph{Uso de Buffer a la entrada (Seguidor de Tensión).}
La interconexión de múltiples etapas analógicas, así como la conexión del filtro hacia los bloques de conmutación (como el \textit{Sample and Hold} o la Llave Analógica), presenta un desafío crítico de carga. Para mitigar esto, se ha incorporado un amplificador operacional en configuración de seguidor de tensión (buffer) a la entrada y/o salida de la red de filtrado.

El buffer presenta una impedancia de entrada teóricamente infinita y una impedancia de salida virtualmente nula. Esto cumple una función de \textbf{aislamiento o acoplamiento de impedancias}: 
\begin{itemize}
	\item Evita que la etapa de filtrado demande corriente excesiva de la fuente de señal previa, previniendo caídas de tensión y distorsión armónica.
	\item Otorga la capacidad de manejo de corriente (drive) necesaria para alimentar las etapas posteriores sin que los componentes reactivos del filtro (resistencias y capacitores) "vean" la impedancia de la carga, lo cual alteraría las frecuencias de corte $\omega_0$ y el factor de calidad $Q$ calculados teóricamente.
\end{itemize}
\newpage
\paragraph{Capacitores de Desacople (Bypass).}
En la práctica, las fuentes de alimentación e instrumentación de laboratorio introducen ruido de alta frecuencia, y los propios amplificadores operacionales (TL082) generan picos de corriente transitoria durante los cambios rápidos de señal. 

Para salvaguardar la estabilidad del sistema, se ha añadido un par de capacitores de desacople (típicamente cerámicos de $100\text{ nF}$) conectados lo más físicamente cerca posible de los pines de alimentación ($+V_{cc}$ y $-V_{ee}$) de cada circuito integrado hacia el plano de masa (GND). 
Estos componentes actúan como un cortocircuito para las componentes de interferencia de alta frecuencia, derivando el ruido a tierra. Simultáneamente, funcionan como reservas locales de energía, suministrando los picos de corriente instantáneos que el amplificador requiere sin provocar caídas de tensión en las líneas de alimentación, previniendo así oscilaciones indeseadas en el circuito.

\paragraph{Esquemático Modificado.}
Considerando las adiciones previamente justificadas, el diseño final a implementar en la placa de circuito impreso se detalla en la \textbf{Figura \ref{fig:esquema_modificado}}. Esta topología garantiza tanto la selectividad espectral exigida (Chebyshev de 5to orden) como la robustez eléctrica necesaria para el acondicionamiento de la señal.

\begin{figure}[H]
	\centering
	\includegraphics[width=0.9\textwidth]{Imagenes/3_filtro/Filtro_Completo.pdf} 
	\caption{Esquemático final del filtro incorporando un seguidor de voltaje y capacitores de desacople en las líneas de alimentación.}
	\label{fig:esquema_modificado}
\end{figure}
\newpage
\subsubsection{BONUS: Análisis de Retardo de Grupo Constante y su Aplicabilidad}
La aproximación de Chebyshev utilizada en el diseño del filtro prioriza una alta selectividad en la banda de transición (caída abrupta) a expensas de sacrificar la linealidad de la fase. Esto produce un retardo de grupo ($\tau_g = -d\phi/d\omega$) que varía con la frecuencia, lo que implica que distintas componentes armónicas de una señal compleja sufren desfasajes diferentes al atravesar el sistema.

\paragraph{Aplicaciones sensibles al dominio temporal.}
Existen escenarios donde mantener la integridad morfológica de la señal en el dominio del tiempo es el requisito primario, volviendo inaceptable el uso de un filtro de Chebyshev. Un ejemplo clásico es la instrumentación biomédica, como la adquisición de señales de Electrocardiograma (ECG). 

El complejo QRS de un latido cardíaco está compuesto por una sumatoria precisa de múltiples frecuencias. Si esta señal atravesara nuestro filtro de Chebyshev, la distorsión de fase alteraría la posición relativa de estos armónicos en el tiempo. Como resultado, la onda se deformaría (los picos se ensancharían o se generarían oscilaciones parásitas), lo cual podría derivar en un diagnóstico clínico erróneo, como una falsa arritmia. 

Para estos casos, la síntesis del filtro debe realizarse obligatoriamente utilizando la aproximación de \textbf{Bessel-Thomson}. Esta familia de filtros garantiza un retardo de grupo máximamente plano en la banda de paso. Al tener un retardo constante, todas las componentes de la señal se retrasan exactamente la misma cantidad de tiempo, preservando la forma de la onda original con nula distorsión temporal.

\paragraph{Análisis de viabilidad para el Trabajo Práctico.}
A la luz de los requerimientos de esta práctica, surge la interrogante fundamental: ¿Tendría sentido sintetizar nuestro circuito impreso utilizando un filtro de Bessel en lugar del de Chebyshev?

La respuesta es \textbf{negativa}. En el contexto de los ensayos dictados en la consigna, las aplicaciones principales son el muestreo de tonos senoidales, el análisis de espectros y, de manera destacada, el remuestreo y demodulación de señales de Amplitud Modulada (AM).
\begin{itemize}
	\item \textbf{Prioridad frecuencial:} El objetivo de los filtros en este sistema (tanto el Anti-alias como el Recuperador) es puramente frecuencial: bloquear tajantemente el espectro por encima de la frecuencia de Nyquist para evitar \textit{aliasing} y suprimir las réplicas post-muestreo. 
	\item \textbf{Inmunidad de la información:} En una señal AM, la información útil reside en la variación de amplitud de la envolvente y en la magnitud de sus bandas laterales. La alineación exacta de la fase de las componentes de alta frecuencia no altera la demodulación ni el cálculo de potencia recuperada.
\end{itemize}
\newpage
Para ilustrar esta limitación, la \textbf{Figura \ref{fig:comparativa_filtros}} presenta una simulación comparativa de la respuesta en frecuencia de ambas aproximaciones asumiendo un orden $n=5$ y una $f_p = 22.05\text{ kHz}$. Se observa claramente cómo la topología de Bessel fracasa en alcanzar los $40\text{ dB}$ de atenuación en la banda de rechazo, evidenciando la necesidad de recurrir a la aproximación de Chebyshev para cumplir con la plantilla exigida sin incrementar desmesuradamente la complejidad del hardware.

\begin{figure}[H]
	\centering
	\includegraphics[width=0.8\textwidth]{Imagenes/3_filtro/filtros_comparativa.pdf}
	\caption{Comparativa de respuesta en magnitud y retardo de grupo para filtros pasabajos de Chebyshev y Bessel de 5to orden ($f_p = 22.05\text{ kHz}$).}
	\label{fig:comparativa_filtros}
\end{figure}

\paragraph{Conclusión del Trade-off.}
Si se intentara sintetizar un filtro de Bessel que cumpliera con la misma estricta plantilla de atenuación exigida ($40\text{ dB}$ de rechazo en $f_a = 44.1\text{ kHz}$), la bajísima selectividad inherente a esta aproximación obligaría a elevar el orden del filtro dramáticamente (típicamente a $n \ge 8$). 

Por lo tanto, para las aplicaciones de este trabajo práctico, el uso de Chebyshev es la decisión de diseño óptima. La distorsión de fase es un costo aceptable e irrelevante para el tratamiento de estas señales, a cambio de lograr una excelente supresión espectral con un orden bajo ($n=5$). Implementar un filtro de Bessel aquí aportaría una linealidad de fase innecesaria, penalizando severamente el hardware: exigiría poner en cascada múltiples etapas Sallen-Key adicionales, lo que multiplicaría el consumo, el tamaño del PCB, el ruido térmico y la inestabilidad ante la tolerancia comercial de capacitores y resistencias.